% Projektbeschreibung für Förderanträge (Prototype Fund, Stiftungen, § 123 SGB XI).
% Quellen: README.md, docs/COMPLIANCE.md, docs/funding-outline.md,
% docs/proposal-organisations.md, die Zusammenfassungen der Läufe vom 22.10.2026
% und .orchestration/runs/2026-10-23-dossier/scout-applications.md.
% Markierungen aus dem Scout: V = an der Quelle geprüft, A = archivierte Fassung,
% S = Sekundärquelle, TV = zu prüfen. A, S und TV stehen hier nur mit \todo{}.
% Bauen aus docs/dossier/antrag/:
%   latexmk -pdf -interaction=nonstopmode -outdir=../build projektbeschreibung.tex
\documentclass[11pt,a4paper,parskip=half-]{scrartcl}
% Shared preamble for the Anspruch dossier and its companion documents.
% Load after \documentclass with \input{preamble} (from docs/dossier/)
% or % Shared preamble for the Anspruch dossier and its companion documents.
% Load after \documentclass with \input{preamble} (from docs/dossier/)
% or % Shared preamble for the Anspruch dossier and its companion documents.
% Load after \documentclass with \input{preamble} (from docs/dossier/)
% or \input{../preamble} (from docs/dossier/antrag/ or docs/dossier/partner/).
% Each document sets its own \documentclass and its own pdftitle via
% \hypersetup{pdftitle={...}}. Only packages that ship with MiKTeX.

\usepackage[T1]{fontenc}
\usepackage[utf8]{inputenc}
\usepackage{lmodern}
\usepackage[english,ngerman]{babel} % last option = main language: ngerman
\usepackage[a4paper,margin=25mm]{geometry}
\usepackage{booktabs}
\usepackage{longtable}
\usepackage{array}
\usepackage{enumitem}
\usepackage{xcolor}
\usepackage{microtype}
\usepackage[autostyle=true,german=quotes]{csquotes}
\usepackage{marginnote}
\usepackage[hidelinks,unicode,pdflang={de-DE}]{hyperref}

% ---------------------------------------------------------------------------
% Colours
\definecolor{blankbg}{RGB}{255,236,153}   % highlighted placeholder
\definecolor{todofg}{RGB}{170,40,40}      % "to confirm" margin note

% \blank{Text}: a highlighted placeholder "[Text]" for values the sender
% fills in before sending. Usable in running text and in table cells.
\newcommand{\blank}[1]{\begingroup\setlength{\fboxsep}{1pt}%
  \colorbox{blankbg}{\strut[#1]}\endgroup}

% \todo{Text}: a "zu bestätigen" note in the margin. Works in running text,
% lists and longtable cells (via marginnote, not \marginpar).
% Defined with \providecommand so that it does not clash if a document
% loads todonotes on its own (todonotes then wins).
\providecommand{\todo}[1]{\marginnote{\raggedright\scriptsize
  \textcolor{todofg}{\textbf{Zu bestätigen:} #1}}}
\setlength{\marginparwidth}{20mm}

% \draftmark: header text for every page of a draft.
\newcommand{\draftmark}{Entwurf, 23.10.2026}

% Draft header for scrartcl / scrreprt (KOMA) and article / report.
% Letters (scrlttr2, letter) place \draftmark themselves.
\makeatletter
\@ifclassloaded{scrartcl}{\def\anspruch@koma{}}{}
\@ifclassloaded{scrreprt}{\def\anspruch@koma{}}{}
\@ifclassloaded{article}{\def\anspruch@std{}}{}
\@ifclassloaded{report}{\def\anspruch@std{}}{}
\ifdefined\anspruch@koma
  \usepackage[headsepline]{scrlayer-scrpage}
  \clearpairofpagestyles
  \ohead{\small\draftmark}
  \ihead{\small\headmark}
  \cfoot*{\pagemark}
  \automark[section]{section}
  \pagestyle{scrheadings}
\fi
\ifdefined\anspruch@std
  \usepackage{fancyhdr}
  \pagestyle{fancy}
  \fancyhf{}
  \fancyhead[R]{\small\draftmark}
  \fancyhead[L]{\small\nouppercase{\leftmark}}
  \fancyfoot[C]{\thepage}
  \fancypagestyle{plain}{\fancyhf{}\fancyhead[R]{\small\draftmark}%
    \fancyfoot[C]{\thepage}\renewcommand{\headrulewidth}{0pt}}
\fi
\makeatother

% Lists: compact, as in a German Antrag.
\setlist{itemsep=0.2em,topsep=0.4em}

% Tables: a bit more air, ragged-right paragraph columns.
\renewcommand{\arraystretch}{1.15}
\newcolumntype{L}[1]{>{\raggedright\arraybackslash}p{#1}}
 (from docs/dossier/antrag/ or docs/dossier/partner/).
% Each document sets its own \documentclass and its own pdftitle via
% \hypersetup{pdftitle={...}}. Only packages that ship with MiKTeX.

\usepackage[T1]{fontenc}
\usepackage[utf8]{inputenc}
\usepackage{lmodern}
\usepackage[english,ngerman]{babel} % last option = main language: ngerman
\usepackage[a4paper,margin=25mm]{geometry}
\usepackage{booktabs}
\usepackage{longtable}
\usepackage{array}
\usepackage{enumitem}
\usepackage{xcolor}
\usepackage{microtype}
\usepackage[autostyle=true,german=quotes]{csquotes}
\usepackage{marginnote}
\usepackage[hidelinks,unicode,pdflang={de-DE}]{hyperref}

% ---------------------------------------------------------------------------
% Colours
\definecolor{blankbg}{RGB}{255,236,153}   % highlighted placeholder
\definecolor{todofg}{RGB}{170,40,40}      % "to confirm" margin note

% \blank{Text}: a highlighted placeholder "[Text]" for values the sender
% fills in before sending. Usable in running text and in table cells.
\newcommand{\blank}[1]{\begingroup\setlength{\fboxsep}{1pt}%
  \colorbox{blankbg}{\strut[#1]}\endgroup}

% \todo{Text}: a "zu bestätigen" note in the margin. Works in running text,
% lists and longtable cells (via marginnote, not \marginpar).
% Defined with \providecommand so that it does not clash if a document
% loads todonotes on its own (todonotes then wins).
\providecommand{\todo}[1]{\marginnote{\raggedright\scriptsize
  \textcolor{todofg}{\textbf{Zu bestätigen:} #1}}}
\setlength{\marginparwidth}{20mm}

% \draftmark: header text for every page of a draft.
\newcommand{\draftmark}{Entwurf, 23.10.2026}

% Draft header for scrartcl / scrreprt (KOMA) and article / report.
% Letters (scrlttr2, letter) place \draftmark themselves.
\makeatletter
\@ifclassloaded{scrartcl}{\def\anspruch@koma{}}{}
\@ifclassloaded{scrreprt}{\def\anspruch@koma{}}{}
\@ifclassloaded{article}{\def\anspruch@std{}}{}
\@ifclassloaded{report}{\def\anspruch@std{}}{}
\ifdefined\anspruch@koma
  \usepackage[headsepline]{scrlayer-scrpage}
  \clearpairofpagestyles
  \ohead{\small\draftmark}
  \ihead{\small\headmark}
  \cfoot*{\pagemark}
  \automark[section]{section}
  \pagestyle{scrheadings}
\fi
\ifdefined\anspruch@std
  \usepackage{fancyhdr}
  \pagestyle{fancy}
  \fancyhf{}
  \fancyhead[R]{\small\draftmark}
  \fancyhead[L]{\small\nouppercase{\leftmark}}
  \fancyfoot[C]{\thepage}
  \fancypagestyle{plain}{\fancyhf{}\fancyhead[R]{\small\draftmark}%
    \fancyfoot[C]{\thepage}\renewcommand{\headrulewidth}{0pt}}
\fi
\makeatother

% Lists: compact, as in a German Antrag.
\setlist{itemsep=0.2em,topsep=0.4em}

% Tables: a bit more air, ragged-right paragraph columns.
\renewcommand{\arraystretch}{1.15}
\newcolumntype{L}[1]{>{\raggedright\arraybackslash}p{#1}}
 (from docs/dossier/antrag/ or docs/dossier/partner/).
% Each document sets its own \documentclass and its own pdftitle via
% \hypersetup{pdftitle={...}}. Only packages that ship with MiKTeX.

\usepackage[T1]{fontenc}
\usepackage[utf8]{inputenc}
\usepackage{lmodern}
\usepackage[english,ngerman]{babel} % last option = main language: ngerman
\usepackage[a4paper,margin=25mm]{geometry}
\usepackage{booktabs}
\usepackage{longtable}
\usepackage{array}
\usepackage{enumitem}
\usepackage{xcolor}
\usepackage{microtype}
\usepackage[autostyle=true,german=quotes]{csquotes}
\usepackage{marginnote}
\usepackage[hidelinks,unicode,pdflang={de-DE}]{hyperref}

% ---------------------------------------------------------------------------
% Colours
\definecolor{blankbg}{RGB}{255,236,153}   % highlighted placeholder
\definecolor{todofg}{RGB}{170,40,40}      % "to confirm" margin note

% \blank{Text}: a highlighted placeholder "[Text]" for values the sender
% fills in before sending. Usable in running text and in table cells.
\newcommand{\blank}[1]{\begingroup\setlength{\fboxsep}{1pt}%
  \colorbox{blankbg}{\strut[#1]}\endgroup}

% \todo{Text}: a "zu bestätigen" note in the margin. Works in running text,
% lists and longtable cells (via marginnote, not \marginpar).
% Defined with \providecommand so that it does not clash if a document
% loads todonotes on its own (todonotes then wins).
\providecommand{\todo}[1]{\marginnote{\raggedright\scriptsize
  \textcolor{todofg}{\textbf{Zu bestätigen:} #1}}}
\setlength{\marginparwidth}{20mm}

% \draftmark: header text for every page of a draft.
\newcommand{\draftmark}{Entwurf, 23.10.2026}

% Draft header for scrartcl / scrreprt (KOMA) and article / report.
% Letters (scrlttr2, letter) place \draftmark themselves.
\makeatletter
\@ifclassloaded{scrartcl}{\def\anspruch@koma{}}{}
\@ifclassloaded{scrreprt}{\def\anspruch@koma{}}{}
\@ifclassloaded{article}{\def\anspruch@std{}}{}
\@ifclassloaded{report}{\def\anspruch@std{}}{}
\ifdefined\anspruch@koma
  \usepackage[headsepline]{scrlayer-scrpage}
  \clearpairofpagestyles
  \ohead{\small\draftmark}
  \ihead{\small\headmark}
  \cfoot*{\pagemark}
  \automark[section]{section}
  \pagestyle{scrheadings}
\fi
\ifdefined\anspruch@std
  \usepackage{fancyhdr}
  \pagestyle{fancy}
  \fancyhf{}
  \fancyhead[R]{\small\draftmark}
  \fancyhead[L]{\small\nouppercase{\leftmark}}
  \fancyfoot[C]{\thepage}
  \fancypagestyle{plain}{\fancyhf{}\fancyhead[R]{\small\draftmark}%
    \fancyfoot[C]{\thepage}\renewcommand{\headrulewidth}{0pt}}
\fi
\makeatother

% Lists: compact, as in a German Antrag.
\setlist{itemsep=0.2em,topsep=0.4em}

% Tables: a bit more air, ragged-right paragraph columns.
\renewcommand{\arraystretch}{1.15}
\newcolumntype{L}[1]{>{\raggedright\arraybackslash}p{#1}}

\hypersetup{pdftitle={Anspruch: Projektbeschreibung}}


\begin{document}
{\LARGE\sffamily\bfseries Anspruch: Projektbeschreibung\par}
\blank{Name der Antragstellerin oder des Antragstellers}\par

\section*{Vor dem Versand ausfüllen}
Jeder gelb markierte Platzhalter muss vor dem Versand ersetzt werden.
Randnotizen in Rot sind offene Prüfpunkte; sie werden vor dem Versand geklärt und gelöscht.
\begin{quote}\raggedright
\blank{Name der Antragstellerin oder des Antragstellers} \textperiodcentered{}
\blank{Adresse der Seite} \textperiodcentered{}
\blank{Programm und Spur} \textperiodcentered{}
\blank{Vergleich mit ähnlichen Angeboten} \textperiodcentered{}
\blank{Ergebnis AP1} bis \blank{Ergebnis AP5} \textperiodcentered{}
\blank{Aufwand AP1} bis \blank{Aufwand AP5} \textperiodcentered{}
\blank{Arbeitsstunden gesamt} \textperiodcentered{}
\blank{Anzahl Pilotpartner} \textperiodcentered{}
\blank{Pilotpartner} \textperiodcentered{}
\blank{Lizenzentscheidung} \textperiodcentered{}
\blank{Team} \textperiodcentered{}
\blank{Erfahrung, bis zu drei Projekte} \textperiodcentered{}
\blank{Land} \textperiodcentered{}
\blank{Kommune oder Träger}
\end{quote}
Die Gliederung folgt dem \enquote{Kleinen Leitfaden für Förderanfragen an Stiftungen}
des Initiativkreises Hamburger Stiftungen, veröffentlicht beim Bundesverband Deutscher
Stiftungen (\href{https://www.stiftungen.org/fileadmin/stiftungen_org/Verband/Was_wir_tun/Veranstaltungen/AK-Foerderstiftungen/Leitfaden-fuer-Foerderanfragen-an-Stiftungen.pdf}{Leitfaden}).
Die Randnotizen \enquote{Formularfeld} nennen das passende Feld im Formular des Prototype Fund
und seine Längengrenze. Felder und Grenzen stammen aus dem Bewerbungshandbuch für Jahrgang 02
(Stand September 2025, archivierte Fassung). Das Formular für Jahrgang 03 ist nur nach Anmeldung
sichtbar; die Grenzen sind dort zu prüfen.

% ---------------------------------------------------------------------------
\section{Kurzfassung}
\todo{Formularfeld 2 \enquote{Projektbeschreibung}, max.\ 100 Wörter (J02, Archiv)}%
Anspruch ist eine Webseite für Familien, die einen Angehörigen pflegen.
Die Familie beantwortet Fragen zum Alltag der pflegebedürftigen Person.
Die Seite schätzt den Pflegegrad, den die Begutachtung wahrscheinlich ergibt.
Sie vergleicht die Leistungen, die daraus folgen, mit dem, was der Haushalt schon bekommt.
Sie nennt die nächsten Schritte in Reihenfolge, jeweils mit der gesetzlichen Grundlage,
und erstellt ein PDF auf Deutsch oder Englisch.
In der Regel genügen 12 bis 20 Fragen statt der 64 Kriterien des amtlichen Instruments.
Die Oberfläche gibt es in sechs Sprachen.
Alles läuft im Browser: Es gibt kein Konto und keinen Server, der Antworten speichert.
Die geförderte Phase macht die Seite bereit für den Einsatz in Beratungsstellen:
muttersprachliche Prüfung, rechtliche Prüfung, Pilote mit Partnern, Anpassung an eine
mögliche Pflegereform und eine Auswertung in Summenzahlen.

% ---------------------------------------------------------------------------
\section{Ausgangslage und Problem}
\todo{Formularfeld 3 \enquote{Gesellschaftliche Herausforderung}, max.\ 175 Wörter (J02, Archiv). Quelle zum nicht abgerufenen Entlastungsbetrag ergänzen.}%
Leistungen der Pflegeversicherung werden nicht von selbst gezahlt.
Sie müssen beantragt werden, oft einzeln, und in einer Sprache, die viele Familien nicht sprechen.
Auch bewilligtes Geld wird oft nicht abgerufen, etwa der Entlastungsbetrag.
Am meisten verlieren dabei die sehr alten Menschen, die erschöpften Angehörigen und Familien,
die den Schriftverkehr in einer zweiten Sprache führen.

Das Gesetz sieht Beratung und Unterstützung vor:
\begin{itemize}
  \item \href{https://www.gesetze-im-internet.de/sgb_11/__7a.html}{§ 7a SGB XI}:
        Anspruch auf individuelle Pflegeberatung; die Pflegekasse benennt unverzüglich
        eine Beraterin oder einen Berater.
  \item \href{https://www.gesetze-im-internet.de/sgb_11/__7b.html}{§ 7b SGB XI}:
        Nach einem Erstantrag bietet die Pflegekasse innerhalb von zwei Wochen einen
        Beratungstermin an oder stellt einen Beratungsgutschein aus.
  \item \href{https://www.gesetze-im-internet.de/sgb_11/__7c.html}{§ 7c SGB XI}:
        Pflegestützpunkte, wenn das Land sie bestimmt, mit unabhängiger Beratung und Koordination.
  \item \href{https://www.gesetze-im-internet.de/sgb_11/__45.html}{§ 45 SGB XI}:
        Unentgeltliche Pflegekurse für Angehörige, auf Wunsch auch zu Hause.
\end{itemize}
Diese Rechte nützen nur, wenn die Familie vorbereitet in die Beratung und in die Begutachtung geht.
Sie muss wissen, welche Leistungen in Frage kommen und was sie gefragt wird.
Hier setzt Anspruch an.

% ---------------------------------------------------------------------------
\section{Zielgruppe}
\todo{Formularfeld 8 \enquote{Zielgruppe}, max.\ 100 Wörter (J02, Archiv)}%
\begin{itemize}
  \item \textbf{Familien}, die einen Angehörigen zu Hause pflegen: vor dem Erstantrag,
        vor der Begutachtung oder wenn der Bedarf gestiegen ist.
  \item \textbf{Pflegende, die Deutsch als zweite Sprache sprechen.}
        Die Oberfläche gibt es auf Deutsch, Englisch, Türkisch, Russisch, Polnisch und Arabisch;
        Arabisch wird durchgehend von rechts nach links dargestellt.
  \item \textbf{Beraterinnen und Berater} in Pflegeberatung, Pflegestützpunkten und
        Beratungsstellen. Sie nutzen die Seite im Gespräch oder geben Familien den Link mit.
\end{itemize}

% ---------------------------------------------------------------------------
\section{Ziel und Nutzen}
Ziel ist, dass Familien vorbereitet zur Pflegekasse, zur Beratung und zur Begutachtung gehen.
\begin{itemize}
  \item Die Familie sieht, welcher Pflegegrad wahrscheinlich ist und welche Leistungen daraus folgen.
  \item Sie sieht die Schritte in Reihenfolge, jeweils mit Paragraf.
  \item Sie nimmt ein PDF mit ihren eigenen Antworten mit. Die Begutachtung fragt dasselbe;
        mit den Antworten auf Papier muss sie sich im Gespräch nicht an alles erinnern.
  \item Die Seite sagt jeder Familie, dass die Pflegeberatung nach § 7a SGB XI kostenlos ist und
        ihr zusteht. Sie schickt Familien zur Beratung, nicht von ihr weg.
  \item Beratende bekommen ein Werkzeug, das im Gespräch in sechs Sprachen funktioniert.
\end{itemize}
Anspruch ist eine Schätzung aus den Angaben der Familie.
Es ist keine Begutachtung durch den Medizinischen Dienst und keine Rechtsberatung.

% ---------------------------------------------------------------------------
\section{Lösung}

\subsection{Was heute besteht}
\todo{Formularfeld 5 \enquote{Aktueller Stand}, max.\ 100 Wörter (J02, Archiv)}%
Die Seite ist gebaut und auf Vercel bereitgestellt (\blank{Adresse der Seite}).
\begin{itemize}
  \item \textbf{Kurze Befragung.} In der Regel 12 bis 20 Fragen statt 64.
        Die Befragung endet, wenn weitere Antworten das Ergebnis nicht mehr ändern können.
        Steht es noch nicht fest, zeigt die Seite einen Bereich und nie einen Mittelwert.
        Der Bericht nennt, wie viele Fragen gestellt wurden und warum die übrigen entfallen.
  \item \textbf{Beträge mit Quelle.} Jeder Betrag und jede Schwelle nennt die Vorschrift
        und das Datum der letzten Prüfung. Alle Beträge wurden am 22.10.2026 gegen das geltende
        Recht geprüft. Ein Betrag ohne Quelle wird nicht angezeigt.
  \item \textbf{Vorsichtige Summen.} Einmalige Zuschüsse zählen nicht zur Monatssumme.
        Leistungen mit Bedingungen erscheinen als \enquote{prüfen, ob das zutrifft}.
        Pflegegeld und Pflegesachleistung werden nicht addiert.
  \item \textbf{PDF auf dem Gerät.} Schätzung, Beträge, Schritte mit Paragraf und die gegebenen
        Antworten, auf Deutsch oder Englisch.
  \item \textbf{Sechs Sprachen} in der Oberfläche. Die Fragen stehen auf Deutsch oder Englisch,
        weil sie dem amtlichen Instrument folgen.
  \item \textbf{Barrierefreiheit.} Einstellungen für Schrift, Farbe, Sprache und Tempo auf der
        ersten Seite; vier Farbschemata mit Kontrast nach WCAG AAA; Vorlesen; Antworten per Sprache
        in Browsern, die Spracherkennung anbieten (Chrome und Edge).
  \item \textbf{Zählung und Partnerlinks.} Ohne Einstellung sendet die Seite nichts.
        Eingeschaltet zählt sie Besuche und fünf Schritte ohne Cookies.
        Partner erhalten einen eigenen Link mit Kürzel.
  \item \textbf{Impressum und Datenschutzerklärung} als Seiten, mit markierten Lücken für die
        Angaben zum Betreiber.
\end{itemize}

\subsection{Technische Umsetzung}
\todo{Formularfeld 4 \enquote{Technische Umsetzung}, max.\ 175 Wörter (J02, Archiv)}%
Anspruch ist eine Web-Anwendung mit Next.js. Schätzung, Speicherung und PDF entstehen im Browser.
\begin{itemize}
  \item \textbf{Regelwerk.} Beträge, Schwellen und die Umrechnungstabellen des
        Begutachtungsinstruments liegen in \texttt{lib/rules/}. Jeder Wert trägt eine Quellen-ID
        mit Vorschrift und Prüfdatum (\texttt{lib/rules/sources.ts}).
  \item \textbf{Abbruchregel.} Jede offene Frage gilt als Spanne. Die Seite rechnet zweimal,
        einmal mit der günstigsten und einmal mit der ungünstigsten Annahme. Stimmen beide
        Grenzen überein, kann keine weitere Antwort den Pflegegrad ändern. Das gilt, weil die
        Bewertung in jedem Modul monoton ist (\texttt{lib/intake/adaptive.ts}).
  \item \textbf{PDF.} Erstellt mit \texttt{jspdf} im Browser; kein Upload, kein Netzwerkaufruf.
  \item \textbf{Speicherung.} Antworten liegen nur im Browser und lassen sich mit einem Tipp löschen.
        Gespeichert wird erst, wenn etwas beantwortet ist.
  \item \textbf{Zählung.} Aus, solange beim Bauen keine Variablen gesetzt sind. Eingeschaltet über
        Plausible, ohne Cookies. Eine Positivliste im Code verhindert, dass Antworten, Pflegegrad,
        Punkte oder Freitext gesendet werden; Tests sichern sie ab.
  \item \textbf{Partnerlinks.} Das Kürzel aus \texttt{?via=} wird geprüft, aus der Adresszeile
        entfernt und nur für den Besuch im Speicher gehalten, nie auf dem Gerät.
\end{itemize}

\subsection{Was neu ist}
\todo{Formularfeld 7 \enquote{Innovation / ähnliche Ansätze}, max.\ 100 Wörter (J02, Archiv)}%
Die Seite verkürzt die Befragung, ohne zu raten: Sie hört erst auf, wenn das Ergebnis feststeht,
und zeigt sonst einen Bereich. Sie braucht keinen Server für Antworten.
Abgrenzung zu ähnlichen Angeboten: \blank{Vergleich mit ähnlichen Angeboten}.

\subsection{Was die geförderte Phase ergänzt}
Es fehlen die muttersprachliche Prüfung, die rechtliche Prüfung, Erfahrung mit Partnern, eine Auswertung
und, falls die Pflegereform kommt, die Anpassung der Beträge. Diese Punkte sind die Arbeitspakete in
Abschnitt~\ref{sec:ap}.

% ---------------------------------------------------------------------------
\section{Vorgehen und Arbeitspakete}\label{sec:ap}
\todo{Formularfeld 9 \enquote{Meilensteine}, max.\ 100 Wörter; Feld 12 \enquote{Arbeitsstunden}: eine Zahl (J02, Archiv)}%
Zeiten und Meilensteine stehen im Zeitplan, Kosten im Finanzplan.
Arbeitsstunden gesamt: \blank{Arbeitsstunden gesamt}.

\subsection*{AP1: Muttersprachliche Prüfung}
Die Oberfläche auf Türkisch, Russisch, Polnisch und Arabisch ist bisher nur von einem
Sprachmodell gegen den deutschen Text geprüft. Je eine Person pro Sprache, die das deutsche
Pflegesystem kennt, liest zuerst die Texte mit gesetzlichem Inhalt
(\texttt{privacyBody}, \texttt{bescheidHint}, \texttt{m5CapBody}, \texttt{exclusiveNote},
\texttt{disclaimerExtra}), dann die neueren Texte
(\texttt{m5SkippedTitle}, \texttt{situationStep}, \texttt{pointsRangeTitle},
\texttt{gradeNoneNominative} und die Hinweise zur Zählung) und schließlich den Rest.
Dabei werden drei offene Fragen entschieden: Polnisch nutzt die Du-Form, wo das Deutsche
\enquote{Sie} sagt; die arabischen Schaltflächen mischen Befehlsform und Verbalsubstantiv;
und ob Pflegegeld, Pflegesachleistung und Pflegekasse beim ersten Vorkommen erklärt werden.
\par\textbf{Ergebnis:} \blank{Ergebnis AP1}\quad \textbf{Aufwand:} \blank{Aufwand AP1}

\subsection*{AP2: Rechtliche und fachliche Prüfung}
Aus \texttt{docs/COMPLIANCE.md}, vor dem Einschalten der Zählung:
\begin{itemize}
  \item Angaben zum Betreiber im Impressum und in der Datenschutzerklärung; zuständige Aufsichtsbehörde.
  \item Wahl des Vercel-Tarifs (Hobby ist für nicht kommerzielle Nutzung) und des Plausible-Tarifs.
  \item Auftragsverarbeitungsverträge mit Vercel und Plausible bestätigen und ablegen.
  \item Kurze rechtliche Prüfung zur Orientierungshilfe Digitale Dienste der DSK
        (Rn.~87--90 zur Reichweitenmessung, Rn.~68/69, Rn.~80 zur Speicherung im Browser)
        und zum Band der monatlichen Lücke in der Zählung.
  \item Verzeichnis der Verarbeitungstätigkeiten mit drei Einträgen; Position zu BFSG und BITV~2.0.
  \item Klären, wie Chrome und Edge die Spracherkennung verarbeiten.
  \item Lizenz entscheiden und den Weg der Doppellizenz anwaltlich bestätigen lassen
        (Abschnitt~\ref{sec:oss}).
\end{itemize}
Fachliche Textfragen aus der Prüfung vom 22.10.2026: ob die Pflegekasse ein englisches PDF annimmt;
ob der Hinweis zu Pflegegeld und Pflegesachleistung die Kombination nach § 38 SGB XI
richtig wiedergibt; ein Hinweistext mit einer Zahl ohne Quellenangabe (\texttt{estimateHigherBody});
die erneute Prüfung der Begutachtungs-Richtlinien, damit das angezeigte Datum des Regelwerks
aktuell wird.
\par\textbf{Ergebnis:} \blank{Ergebnis AP2}\quad \textbf{Aufwand:} \blank{Aufwand AP2}

\subsection*{AP3: Pilote mit Partnern}
Nach dem Vorschlag an Organisationen (\texttt{docs/proposal-organisations.md}):
drei Monate je Partner; zwei Beraterinnen oder Berater, die die Seite im Gespräch nutzen oder
den Link mitgeben; ein Rückmeldegespräch pro Monat; eine Strichliste der Beratenden ohne Namen.
Jeder Partner erhält einen eigenen Link. Geplant sind \blank{Anzahl Pilotpartner} Partner:
\blank{Pilotpartner}. Die Pilote beginnen erst, wenn AP1 und AP2 abgeschlossen sind.
\par\textbf{Ergebnis:} \blank{Ergebnis AP3}\quad \textbf{Aufwand:} \blank{Aufwand AP3}

\subsection*{AP4: Anpassung an das Pflegeneuordnungsgesetz, falls beschlossen}
Das Kabinett hat den Entwurf des Pflegeneuordnungsgesetzes (PNOG) am 30.09.2026 beschlossen.
Die erste Lesung im Bundestag ist für den 05.11.2026 angesetzt. Das Gesetz ist noch nicht beschlossen
und kann sich ändern. Wird es wie entworfen beschlossen, ändern sich die meisten Beträge und die
Schwellen der Pflegegrade ab dem 01.01.2027. Dann müssen Regelwerk, Quellen-IDs, Tests, Texte in
sechs Sprachen und das PDF angepasst werden; neue Texte gehen durch AP1.
Wird es nicht beschlossen, bleibt die Prüfung aller Beträge vor jeder Veröffentlichung.
\par\textbf{Ergebnis:} \blank{Ergebnis AP4}\quad \textbf{Aufwand:} \blank{Aufwand AP4}

\subsection*{AP5: Auswertung}
Monatlich und je Partnerlink die Summenzahlen aus Abschnitt~\ref{sec:wirkung},
dazu die Strichlisten der Beratenden und die Notizen aus den Rückmeldegesprächen.
Am Ende steht ein Bericht ohne personenbezogene Daten.
\par\textbf{Ergebnis:} \blank{Ergebnis AP5}\quad \textbf{Aufwand:} \blank{Aufwand AP5}

% ---------------------------------------------------------------------------
\section{Wirkung und Indikatoren}\label{sec:wirkung}
Die Zählung erfasst Besuche und fünf Schritte. Jeder Schritt trägt die Sprache der Oberfläche,
das Tempo (eine Frage oder Liste) und das Partnerkürzel.

\begingroup\small
\begin{longtable}{@{}L{0.24\linewidth}L{0.30\linewidth}L{0.38\linewidth}@{}}
\toprule
\textbf{Schritt} & \textbf{Wann} & \textbf{Zusätzliche Angaben} \\
\midrule
\endhead
\texttt{intake\_started} & Start auf der ersten Seite & keine \\
\texttt{result\_shown} & Der Bericht erscheint &
  Fragen: bis 10, bis 20, bis 40, über 40; Ergebnis: keines, Bereich, Pflegegrad (ob, nicht welcher);
  monatliche Lücke: 0, unter 100, unter 300, unter 600, ab 600 EUR \\
\texttt{questions\_refined} & Weitere Fragen werden beantwortet & keine \\
\texttt{pdf\_downloaded} & Das PDF ist erstellt & Sprache des PDF: Deutsch oder Englisch \\
\texttt{intake\_erased} & Die Antworten werden gelöscht & keine \\
\bottomrule
\end{longtable}
\endgroup

Dazu kommen die Strichlisten der Beratenden: wie oft die Seite genutzt wurde und wie oft danach
ein Antrag, ein Antrag auf Höherstufung oder ein Widerspruch folgte.
Indikatoren sind diese Zahlen und der Anteil der begonnenen Befragungen, die ein Ergebnis erreichen.
Ein Ergebnis in Zahlen wird vorab nicht versprochen.
Die Zahlen zeigen Nutzung, nicht den Bescheid der Pflegekasse.

% ---------------------------------------------------------------------------
\section{Datenschutz}
Die Antworten der Familie bleiben auf ihrem Gerät; sie liegen nur in diesem Browser und lassen sich
mit einem Tipp löschen.
Es gibt kein Konto und keinen Server, der Antworten speichert, und das PDF entsteht im Browser.
Ist die Zählung eingeschaltet, zählt ein Anbieter, der seine Server nach eigenen Angaben in der EU
betreibt, Besuche und fünf Schritte ohne Cookies, ohne Antworten und ohne Pflegegrad.
Bei Antworten per Sprache erkennt der Browser die Sprache, nicht die Seite; die rechtliche Prüfung
der Datenschutzerklärung ist Teil von AP2.

% ---------------------------------------------------------------------------
\section{Open Source und Nachhaltigkeit}\label{sec:oss}
Der Quellcode hat noch keine Lizenz. Entscheidung: \blank{Lizenzentscheidung}.
\begin{itemize}
  \item \textbf{Empfehlung:} EUPL-1.2. Das Urheberrecht bleibt in einer Hand, und vor jedem
        Beitrag von außen wird eine Beitragsvereinbarung (CLA) geschlossen. Die EUPL-1.2 ist
        OSI-zertifiziert und wechselseitig; ein Partner kann den Code nicht schließen.
  \item \textbf{Prototype Fund:} Die Ergebnisse sollen unter einer Open-Source-Lizenz öffentlich
        zugänglich sein.\todo{Archiv J02; keine Liste zulässiger Lizenzen, EUPL-1.2 nicht genannt}
        Die Entscheidung fällt vor dem 30.11.2026, wenn die Bewerbung kommt.
  \item \textbf{Doppellizenz:} Mit dem Urheberrecht in einer Hand kann neben der offenen Lizenz
        eine gepflegte Fassung für Organisationen angeboten werden. Das muss anwaltlich bestätigt werden.
  \item \textbf{Hessen:} Ein dort nach § 123 SGB XI gefördertes Ergebnis darf nicht kommerziell
        vertrieben werden und muss Dritten diskriminierungsfrei zur Verfügung stehen. Für ein so
        gefördertes Ergebnis passt der Weg der Doppellizenz nicht.
\end{itemize}
Nach der Förderung fallen Kosten für Hosting und Zählung an, dazu die Prüfung aller Beträge vor
jeder Veröffentlichung und Anpassungen bei Gesetzesänderungen. Die Finanzierung danach ist offen;
Pilotgebühren und eine gepflegte Fassung sind Optionen, keine Zusagen.

% ---------------------------------------------------------------------------
\section{Team}
\todo{Formularfeld 10 \enquote{Team}, max.\ 30 Wörter; Feld 11 \enquote{Erfahrung}, max.\ 3 Projekte, 100 Wörter (J02, Archiv)}%
\blank{Team}\quad Erfahrung: \blank{Erfahrung, bis zu drei Projekte}

% ---------------------------------------------------------------------------
\section{Risiken}
\begingroup\small
\begin{longtable}{@{}L{0.22\linewidth}L{0.31\linewidth}L{0.39\linewidth}@{}}
\toprule
\textbf{Risiko} & \textbf{Folge} & \textbf{Umgang} \\
\midrule
\endhead
PNOG wird beschlossen, geändert oder verschoben & Beträge und Schwellen ab 01.01.2027 falsch; Umfang von AP4 offen. &
  AP4 im Plan und im Budget; Anpassung erst nach Verkündung; bis dahin gilt das geltende Recht. \\
Die rechtliche Prüfung verlangt eine Einwilligung für die Zählung & Ohne Einwilligung keine Summenzahlen. &
  Die Zählung bleibt aus; die Seite sendet dann nichts. AP5 stützt sich auf die Strichlisten. \\
Die Speicherung im Browser über den Besuch hinaus wird beanstandet & Die Speicherung muss angepasst werden. &
  Gespeichert wird erst nach einer Antwort; Löschen mit einem Tipp; Prüfung in AP2. \\
Ungeprüfte Übersetzungen werden geglaubt & Familien handeln auf falscher Grundlage. &
  AP1 vor jedem Pilot; die Fragen selbst stehen auf Deutsch oder Englisch. \\
Partner nutzen die Seite wenig & Wenig Zahlen für AP5. &
  Kleiner Pilot mit zwei Beratenden und monatlichem Gespräch; auch geringe Nutzung ist ein Ergebnis. \\
Die Schätzung weicht von der Begutachtung ab & Die Familie erwartet mehr, als bewilligt wird. &
  Bereich statt Mittelwert; Hinweis, dass es eine Schätzung ist; vorsichtige Summen. \\
\bottomrule
\end{longtable}
\endgroup

% ---------------------------------------------------------------------------
\section{Bezug zum Programm}
Beantragt bei: \blank{Programm und Spur}.

\subsection{Prototype Fund, Jahrgang 03}
\todo{Formularfeld 13 \enquote{Motivation}, max.\ 100 Wörter; Felder 14 bis 16 Second Stage (J02, Archiv)}%
Bewerbungsschluss ist der 30.11.2026 für alle drei Spuren
(\href{https://bewerben.prototypefund.de/}{bewerben.prototypefund.de}). Die Spuren laut Seite:
\begin{itemize}
  \item \enquote{Resilienz}: \enquote{Freie Alternativen zu proprietärer, zentralisierter und
        dominierender Software}.
  \item \enquote{Innovation}: \enquote{Innovative Projekte im Bereich Datensicherheit und
        Software-Infrastruktur}.
  \item \enquote{Up and Coming}: \enquote{für alle unter 26 oder in einem Studium oder in einer
        informatischen Ausbildung}.
\end{itemize}
Vorgeschlagen wird \enquote{Resilienz}: Anspruch erledigt im Browser, was sonst einen Server braucht.
Schätzung, gespeicherte Sitzung und PDF entstehen auf dem Gerät, ohne Konto und ohne Upload.
Ob die Spur passt, wird vor der Bewerbung beim Prototype Fund erfragt.

Gefördert werden laut Handbuch nur Einzelpersonen oder Teams bis vier Personen, als Privatperson,
über 18 Jahre, mit Hauptwohnsitz oder GbR-Sitz in Deutschland.
Eine Second Stage muss schon in der ersten Bewerbung beantragt werden.%
\todo{Archiv J02: Zulassung und Second Stage für J03 bestätigen}

\subsection{Modellvorhaben nach § 123 SGB XI}
Grundlage sind die Empfehlungen des GKV-Spitzenverbandes nach § 123 Abs.~3 SGB XI
(\href{https://www.gkv-spitzenverband.de/media/dokumente/pflegeversicherung/richtlinien__vereinbarungen__formulare/2026-06-01_Pflege_Empfehlungen_nach_123_Abs.3_SGB_XI.pdf}{Empfehlungen}).
Sie nennen die \enquote{Digitale Vernetzung der Pflegeangebote}; ein Rechner für den einzelnen Haushalt
ist nicht genannt. Anspruch wäre deshalb ein Baustein eines Vorhabens, das \blank{Kommune oder Träger}
in \blank{Land} trägt. Nötig sind eine Konzeption, eine wissenschaftliche Begleitung nach § 124 SGB XI
und die hälftige Kofinanzierung durch Land oder Kommune.
Nach Land, soweit geprüft:
\begin{itemize}
  \item \textbf{Niedersachsen:} Antragsberechtigt sind auch natürliche Personen; Frist jeweils
        bis 31.08.\ für einen Start im selben Jahr.
  \item \textbf{Berlin:} Natürliche Personen sind nicht antragsberechtigt; die Frist 21.08.2026 ist
        vorbei.
  \item \textbf{Hessen:} Antrag beim Ministerium mit Stellungnahme des Landkreises oder der
        kreisfreien Stadt; das Ergebnis darf nicht kommerziell vertrieben werden.
  \item \textbf{Schleswig-Holstein:} Richtlinie vom 01.10.2025; Antragsteller und Fristen sind nicht
        angegeben und beim Sozialministerium zu erfragen.
\end{itemize}

\subsection{Aktion Mensch}
Die aktuellen Programme enthalten keines für digitale Teilhabe. Neue Programme sind ab 02.11.2026
angekündigt; die nächste Antragsphase beginnt am 01.01.2027. Einzelpersonen sind nicht förderfähig;
ein Antrag ginge nur über einen gemeinnützigen Träger.

\subsection{Deutsche Fernsehlotterie (Stiftung Deutsches Hilfswerk)}
Organisationsdaten bis 10.12.2026, Bewerbung vom 02.11.\ bis 17.12.2026, Entscheidung im Mai 2027,
frühester Start 01.06.2027. Die Gesamtkosten müssen mindestens 50.000 EUR betragen; je Träger
gibt es eine Bewerbung. Ein Antrag ginge nur über einen gemeinnützigen Träger.%
\todo{\enquote{nur gemeinnützige Träger}: Sekundärquelle; Passung prüfen}

Nicht vorgesehen ist ein Antrag bei der Robert Bosch Stiftung: Sie bearbeitet keine unaufgefordert
eingesandten Förderanfragen.

\end{document}
